\section{方法}

\subsection{研究设计与数据来源}

% Word body block 12
研究包含三部分数据，各自承担不同的评价任务（表 1）。通用纳米材料数据采用 Gul 等论文的公开补充工作簿 [3]，读取细胞存活率工作表，仅纳入存活率、浓度和暴露时间均可解析的记录，得到 3889 条记录、144 个文献组。该数据用于独立基准与迁移预训练，当前清洗数据中没有检出 Se 材料记录，归档分组标识与历史 SeNPs 的 24 个分组也没有交集。

% Word body block 13
历史 SeNPs 数据来自研究组归档的验证工作簿，按现有纳入规则从 623 条原始记录中保留 434 条，排除非研究材料记录、非细胞存活率检测记录及缺少必需建模字段的记录。同一篇研究的数据优先使用符合 DOI 格式的标识分组，其他记录回退到原始文章编号，共形成 24 个文献组，最大组包含 54 条记录。该分析保留理化性质缺失的记录，缺失值在训练折内处理。旧稿的 383 条历史整理数据属于另一分析口径，本研究的主结果均以这 434 条记录及其折外预测为准。

% Word body block 14
新增数据为独立实验工作簿中的351条HUVEC报告值，逐条件冻结模型预测文件另有解盲前封存记录。该数据仅用于实验曲线和描述性比较，不参与历史模型训练、特征筛选或调参，也不将由新增曲线计算的IC50反向输入历史模型。

% Word body block 16: 表格逐单元格转写。
\begin{table}[H]
\centering
\caption{三部分数据的分析范围与用途。记录数不等同于独立材料数或实验重复数。}
\label{tab:data_scope}
\small
\begin{tabularx}{\linewidth}{lrrX}
\toprule
数据 & 记录数 & 分组范围 & 用途 \\
\midrule
通用纳米材料 & 3889 & 144 个文献组 & 通用基准与迁移预训练 \\
历史 SeNPs & 434 & 24 个文献组 & 嵌套分组交叉验证与增量比较 \\
新增 HUVEC 实验 & 351 & 9 种材料、27 条曲线 & 剂量—时间观测、盲法预测评价与表观 IC50 \\
\bottomrule
\end{tabularx}

\end{table}

\subsection{输入特征与折内预处理}

% Word body block 19
预测目标为细胞存活率（\%）。基础配置包含 19 个候选字段，覆盖修饰信息、合成方法、表面电荷、细胞类型与组织、发表年份、粒径、Zeta 电位及其绝对值、暴露时间、浓度、相应的对数变换、细胞接种数，以及 Zeta 电位、接种数、修饰信息和表面电荷的缺失标志。字段名称与转换规则见补充材料，发表年份作为归档背景变量保留，其作用不能解释为材料性质的因果效应。

% Word body block 20
粒径按原始粒径、水中尺寸、培养介质中尺寸的顺序取第一个可用值，Zeta 电位优先采用水中报告值，再回退到培养介质中的报告值。粒径、时间和浓度采用 log(1 + x) 变换，转换不依赖目标值。历史浓度直接使用归档字段的 \ensuremath{\mu}g/mL 数值，不在本轮重新推定或换算。新增实验的原始浓度单位和模型输入换算另行记录，最终解释以单位核验为前提。

% Word body block 21
数值缺失值采用当前训练折的中位数填补，类别缺失值采用当前训练折的众数填补，再进行独热编码。对训练中未出现的类别，编码器忽略未知类别，当前训练折内完全缺失的字段不进入该折模型。以上处理均在内层或外层对应的训练部分拟合，验证部分只接受已拟合的转换，避免用验证数据估计插补值和类别结构。

% Word body block 22
化学增强候选加入元素 Se 含量、稳定性时间、Se 3d 结合能和晶体信息字段，比较基础配置、逐个加入新增字段及联合加入全部字段。当前晶体信息变量记录的是晶体学字段是否有报告，不能当作统一测定的晶相类别。含实测 IC50 的配置单列为响应相关辅助分析，候选包括基础配置、基础加 IC50、单个化学字段及联合配置，不将该配置认定为新材料毒性初筛的主模型。修饰物类别和分子量缺少完整字段，因此当前基础配置也不能视为原计划中四项材料描述符的最小输入模型。

\subsection{模型训练与分组评价}

% Word body block 24
历史主分析采用随机森林（random forest，RF）[6]与XGBoost[7]回归模型。外层采用5折文献分组交叉验证，内层在外层训练数据中采用3折文献分组验证，以内层各折平均绝对误差（mean absolute error，MAE）的平均值选择特征配置和参数。相同文献组不同时出现在对应训练与验证部分，全部434条记录各获得一次外层折外（out-of-fold，OOF）预测。

% Word body block 25
RF候选均使用250棵树，比较不限深度或最大深度为10、最小叶节点样本数为2或4的三种组合。XGBoost比较150或250轮、最大深度2或3的三种组合，学习率均为0.05，行与列采样比例均为0.9，L2L2\hspace{0pt}正则系数为1，采用平方误差目标与直方图树方法。基础随机种子为42，各外层折使用相应折编号偏移后的种子；完整候选表与逐折选择日志保留用于复核。

% Word body block 26
汇总指标包括MAE、均方根误差（root mean square error，RMSE）与决定系数\nobreakspace{}R2，由全部OOF预测计算：

% Word body block 29
其中 yi 和 \ensuremath{\hat{y}}i 分别为存活率报告值与 OOF 预测值，\ensuremath{\bar{y}} 为全部评价记录的平均存活率。MAE 与RMSE 的单位为存活率百分点，R2 无量纲，训练集拟合指标不作为外层泛化结果。

% Word body block 30
增量比较使用同一记录的配对绝对误差，定义 di = |yi \ensuremath{\mathord{-}} \ensuremath{\hat{y}}i,variant \ensuremath{\mathord{-}} |yi \ensuremath{\mathord{-}} \ensuremath{\hat{y}}i,baseline |，负值表示新增配置误差下降。分别报告按样本等权汇总的 \ensuremath{\mathnormal{\Delta}}MAE，以及先在每个文献组内平均 di 、再对 24个组等权平均的差值。配对不确定性采用文献组等权差值的 5000 次有放回 bootstrap，取 2.5\% 与97.5\% 分位数为区间，随机种子为 42。区间与文献组等权估计量配套报告，不能直接视为按样本汇总差值的置信区间；区间整体低于 0 时才支持该比较下的稳定误差降低。

% Word body block 31
通用基准使用固定 500 棵树或提升轮次的 5 折外层文献分组评价，不进行上述内层特征与参数搜索。General 到 SeNP 迁移也采用独立的固定 500 树预算，对比 SeNP-only、General zero-shot和 General\ensuremath{\rightarrow}SeNP 三种策略。RF 迁移保留 250 棵通用数据训练树，再添加 250 棵 SeNPs 训练树，原有树不重新优化；XGBoost 在 250 轮通用预训练后续训 250 轮。迁移预处理仅使用通用数据与当前 SeNPs 训练折，外层 SeNPs 验证记录不参与拟合，这一固定预算对照与嵌套主分析的基线数值分别报告。

% Word body block 32
剂量—时间消融在相同历史记录与 5\ensuremath{\times}3 嵌套分组评价下，比较基础配置与加入 log(1 +C) log(1 + T ) 的配置，其中 C 与 T 分别为归档浓度和时间。该实验检验显式交互项的预测增量，与新增湿实验中观察时间依赖性属于不同问题。

\subsection{新增 HUVEC 实验数据与描述性分析}

% Word body block 34
实验数据包括 CS、Glucose、LET、PAH、PEG、PEI、PVP、Sucrose 和 TW80 九种 SeNPs，细胞模型为人脐静脉内皮细胞（human umbilical vein endothelial cells，HUVEC）。每种材料覆盖 24、48、72 h，浓度报告值为 0、0.25、0.5、1、2、4、8、16、32、64、128、256 和 512，共 9 \ensuremath{\times} 3 \ensuremath{\times} 13 = 351 个条件报告值。每种材料有一组固定报告粒径和 Zeta 电位，当前表格没有这些描述符的重复测定数据。

% Word body block 35
原表浓度单位的含义、是否按元素硒计量及换算基准尚未由实验记录确认，因此本文保留报告浓度尺度，并将模型输入浓度换算视为需核验的分析前提。材料缩写沿用原始记录，除表中明确提供的名称外，不从缩写反推修饰物类别或分子量。【待核实：核对浓度单位、元素硒计量基准，以及 CS、LET、PAH 等材料缩写的完整定义。】

% Word body block 36
每个条件在归档表中只有一个报告值，未提供原始重复孔数量、独立实验次数或汇总方法。本轮保留全部报告值，包括大于\nobreakspace{}100\%的存活率，不裁剪、不插补，也不做额外对照归一化。曲线只连接报告点，不作平滑或单调拟合，不绘制误差条，不开展基于独立重复的显著性检验。【待核实：补充材料制备与表征方法、细胞来源及培养条件、接种密度、毒性检测方法和仪器、空白及对照设置、存活率计算式、技术重复和独立实验次数。】

% Word body block 37
描述性比较采用最高报告浓度512，分别列出9种材料在三个时间点的存活率，并对材料等权平均，以概括高浓度下的时间变化。材料排序限定在HUVEC、72 h、报告浓度512的条件，粒径或Zeta电位与该浓度下存活率的探索性关联采用Spearman秩相关，并列值使用平均秩，每个时间点的样本单位为9种材料，不把351个条件报告值视为351个独立材料。

\subsection{表观半数抑制浓度的估计}

% Word body block 39
对每条材料—时间曲线按浓度排序，在正浓度报告点中查找存活率从高于 50\% 降到不高于 50\% 的相邻括区。设相邻点为 (c\textsubscript{1}, v\textsubscript{1}) 与 (c\textsubscript{2}, v\textsubscript{2})，在 log\textsubscript{1}\textsubscript{0} 浓度尺度上线性插值。

% Word body block 40
log\textsubscript{1}\textsubscript{0}(IC\textsubscript{5}\textsubscript{0}) = log\textsubscript{1}\textsubscript{0}(c\textsubscript{1}) + [(50 \ensuremath{\mathord{-}} v\textsubscript{1})/(v\textsubscript{2} \ensuremath{\mathord{-}} v\textsubscript{1})] \ensuremath{\times} [log\textsubscript{1}\textsubscript{0}(c\textsubscript{2}) \ensuremath{\mathord{-}} log\textsubscript{1}\textsubscript{0}(c\textsubscript{1})]。

% Word body block 41
零浓度对照不参与对数插值，但保留在曲线展示中。存在多个方向相反的 50\% 穿越时不强行给出唯一估计，正浓度范围内始终高于 50\% 的曲线记为 > 512，表示测试范围内没有观测到该阈值，而非外推得到确定数值。估计结果称为表观 IC50，保留相邻点与穿越状态，不计算缺少原始重复数据信息条件下的实验置信区间。模型侧 IC50 也必须满足预测曲线在测试范围内跨越 50\% 的条件，不能仅依据单独的预测 IC50 数值表进行验证。

\subsection{冻结模型的盲法预测与配对误差评价}

% Word body block 43
冻结模型对 9 种新材料、13 个报告浓度和 24、48、72 h 三个时间点生成逐条件预测，共 351 个预测值。封存记录注明封存时间，记录的 SHA-256，当前文件重新计算得到相同哈希。预测表与实验表按材料名称、报告浓度、时间和浓度单位一对一匹配，获得 351 对预测—实测值。

% Word body block 44
以实验存活率为参照，计算 MAE、RMSE、决定系数 R\textsuperscript{2} 和平均偏差。平均偏差定义为预测值减去实测值，正值表示模型高估存活率。另按 24、48、72 h 分层计算 MAE 与 RMSE。由于每个条件仅有一个归档报告值，且未提供原始重复孔数与汇总方法，这些指标描述当前条件表中的配对预测表现，不作为独立实验重复的置信区间。模型输入浓度使用归档文件提供的换算字段，单位含义和计量基准仍需实验记录确认。

\subsection{预测曲线 IAE 的处理边界}

% Word body block 46
任务 24 的归档文件提供 27 条材料—时间曲线的 IAE 汇总值，但当前材料中未找到足以独立复现全部数值的计算脚本和完整公式说明。正式正文暂不将该热图作为经核验的定量证据。若补入正式稿，应明确积分横轴、零浓度点处理、插值方法、积分区间及数值积分规则，并从逐条件预测和实测数据重新计算。

\subsection{预测—实测排序与 IC50 阈值可识别性}

% Word body block 48
材料排序比较限定在 HUVEC、72 h、报告浓度 512 的同一条件下，分别按实测和预测存活率由低到高排序，采用 Spearman 秩相关系数和 Kendall 秩相关系数评价一致性，统计单位为 9 种材料。模型侧 IC50 仅在预测曲线于测试范围内跨越 50\% 存活率时才可由相邻点插值识别。当前 27 条预测曲线的最低预测存活率仍高于 50\%，因此不报告模型侧有限 IC50，也不计算基于这些预测曲线的预测—实测 IC50 配对误差。

\subsection{SHAP 可解释性分析与缺失字段处理}

% Word body block 50
对冻结随机森林模型采用精确 Tree SHAP，设置 tree\_path\_dependent，并使用冻结模型保存的历史训练路径样本计数作为背景。不重新训练模型，也不将新增实验存活率用作 SHAP 输入或背景。粒径组由 size\_final 与 log\_size 的 SHAP 贡献加和得到，Zeta 组由 zeta\_final、abs\_zeta 与 zeta\_missing 的贡献加和得到。Surface\_Charge 类别贡献单独保留，不并入 Zeta 数值组。完整模型的 SHAP 加和校验仍使用全部 116 个编码列，要求 base value 与全部 SHAP 值之和和模型预测之间的最大绝对误差不超过 10\textsuperscript{-}\textsuperscript{6} 个存活率百分点。351 个点表示 9 种材料在不同浓度和时间下的预测条件，不是 351 个独立实验重复。

% Word body block 51
历史数据未提供统一编码的修饰物类别（Modifier Category）和经核实的分子量（MW），新增实验表也缺少这两个字段，因此本研究未将其纳入模型，也未对整列缺失信息进行统计插补或未经核实的人工赋值。现有 Coat/Functional Group 名称字段不等同于标准化 Modifier Category，故当前 SHAP 分析不能评价 Modifier Category 或 MW 的独立贡献。
